\documentclass[12pt, a4paper]{article}

\usepackage[utf8]{inputenc}
\usepackage[T1]{fontenc}
\usepackage[brazil]{babel}
\usepackage{amsmath, amssymb}
\usepackage{geometry}
\usepackage{listings}
\usepackage{xcolor}
\usepackage{enumitem}
\usepackage{mdframed}

\geometry{margin=2.5cm}

% --- Estilo dos blocos de código ---
\lstdefinestyle{cstyle}{
    language=C,
    basicstyle=\ttfamily\small,
    keywordstyle=\color{blue}\bfseries,
    commentstyle=\color{gray}\itshape,
    stringstyle=\color{orange},
    numbers=left,
    numberstyle=\tiny\color{gray},
    stepnumber=1,
    frame=single,
    backgroundcolor=\color{gray!8},
    breaklines=true,
    tabsize=4,
}

\lstdefinestyle{javastyle}{
    language=Java,
    basicstyle=\ttfamily\small,
    keywordstyle=\color{purple}\bfseries,
    commentstyle=\color{gray}\itshape,
    stringstyle=\color{orange},
    numbers=left,
    numberstyle=\tiny\color{gray},
    stepnumber=1,
    frame=single,
    backgroundcolor=\color{purple!5},
    breaklines=true,
    tabsize=4,
}

% --- Caixa de resposta ---
\newenvironment{resposta}{%
    \begin{mdframed}[linecolor=black!30, backgroundcolor=white,
                     innertopmargin=10pt, innerbottommargin=30pt,
                     skipabove=6pt]
    \textbf{Resposta:}\\[4pt]
}{%
    \end{mdframed}
}

% -------------------------------------------------------
\title{\textbf{Lista de Exercícios}\\[6pt]
       \large Transformação de Código em Notação de Somatório}
\author{Lucas Rafael Pessoa do Nascimento}
\date{}

\begin{document}

\maketitle

\begin{center}
    \textit{Para cada exercício, analise o(s) laço(s) do código e expresse o número\\
    total de operações (comparações e/ou atribuições) como um somatório.\\
    Sempre que possível, calcule a forma fechada e indique a complexidade assintótica.}
\end{center}

\bigskip

% =========================================================
\section*{Parte I — Aquecimento: Laços Simples}
% =========================================================

% ---------------------------------------------------------
\subsection*{Exercício 1 — Loop simples em C}

\begin{lstlisting}[style=cstyle]
void imprime(int n) {
    int i;
    for (i = 0; i < n; i++) {   // quantas vezes executa?
        printf("%d\n", i);
    }
}
\end{lstlisting}

\begin{enumerate}[label=(\alph*)]
    \item Escreva o somatório que conta o número de execuções do corpo do laço.
    \item Calcule a forma fechada.
\end{enumerate}

\begin{resposta}
\vspace{3cm}
\end{resposta}

% ---------------------------------------------------------
\subsection*{Exercício 2 — Dois laços aninhados em Java}

\begin{lstlisting}[style=javastyle]
void pares(int n) {
    for (int i = 0; i < n; i++) {
        for (int j = 0; j < n; j++) {   // quantas comparacoes?
            System.out.println(i + ", " + j);
        }
    }
}
\end{lstlisting}

\begin{enumerate}[label=(\alph*)]
    \item Escreva o somatório duplo que representa o número total de iterações.
    \item Simplifique para forma fechada e determine $\Theta(\;)$.
\end{enumerate}

\begin{resposta}
\vspace{3cm}
\end{resposta}

% ---------------------------------------------------------
\subsection*{Exercício 3 — Laço triangular em C}

\begin{lstlisting}[style=cstyle]
void triangular(int n) {
    int i, j;
    for (i = 0; i < n; i++) {
        for (j = 0; j < i; j++) {   // limite interno depende de i
            printf("(%d,%d)\n", i, j);
        }
    }
}
\end{lstlisting}

\begin{enumerate}[label=(\alph*)]
    \item Por que este caso é diferente do Exercício 2?
    \item Escreva o somatório e calcule a forma fechada.
    \item Compare o resultado com $\dfrac{n(n-1)}{2}$ — o que você observa?
\end{enumerate}

\begin{resposta}
\vspace{3.5cm}
\end{resposta}

\newpage
% =========================================================
\section*{Parte II — Algoritmos de Ordenação}
% =========================================================

% ---------------------------------------------------------
\subsection*{Exercício 4 — Selection Sort em C}

\begin{lstlisting}[style=cstyle]
void selectionSort(int v[], int n) {
    int i, j, min, tmp;
    for (i = 0; i < n - 1; i++) {
        min = i;
        for (j = i + 1; j < n; j++) {   // comparacao: v[j] < v[min]
            if (v[j] < v[min])
                min = j;
        }
        tmp = v[min]; v[min] = v[i]; v[i] = tmp;
    }
}
\end{lstlisting}

\begin{enumerate}[label=(\alph*)]
    \item Identifique a operação dominante (comparação na linha 6).
    \item Escreva o somatório duplo que conta o número exato de comparações.
    \item Calcule a forma fechada e determine $\Theta(\;)$.
    \item O número de \textit{trocas} é sempre o mesmo? Expresse-o como somatório.
\end{enumerate}

\begin{resposta}
\vspace{4cm}
\end{resposta}

% ---------------------------------------------------------
\subsection*{Exercício 5 — Bubble Sort em Java}

\begin{lstlisting}[style=javastyle]
void bubbleSort(int[] v) {
    int n = v.length;
    for (int i = 0; i < n - 1; i++) {
        for (int j = 0; j < n - 1 - i; j++) {  // comparacao
            if (v[j] > v[j + 1]) {
                int tmp = v[j];
                v[j] = v[j + 1];
                v[j + 1] = tmp;
            }
        }
    }
}
\end{lstlisting}

\begin{enumerate}[label=(\alph*)]
    \item O limite do laço interno muda a cada iteração do externo. Escreva o somatório que reflete isso.
    \item Calcule a forma fechada. O resultado é igual ao do Selection Sort?
    \item No \textbf{melhor caso} (vetor já ordenado), quantas comparações são feitas? E no \textbf{pior caso}? Expresse ambos como somatórios.
\end{enumerate}

\begin{resposta}
\vspace{4cm}
\end{resposta}

% ---------------------------------------------------------
\subsection*{Exercício 6 — Insertion Sort em C}

\begin{lstlisting}[style=cstyle]
void insertionSort(int v[], int n) {
    int i, j, chave;
    for (i = 1; i < n; i++) {
        chave = v[i];
        j = i - 1;
        while (j >= 0 && v[j] > chave) {   // comparacao
            v[j + 1] = v[j];               // deslocamento
            j--;
        }
        v[j + 1] = chave;
    }
}
\end{lstlisting}

\begin{enumerate}[label=(\alph*)]
    \item O número de iterações do \texttt{while} depende dos dados. Escreva o somatório para o \textbf{pior caso} (vetor em ordem decrescente).
    \item Escreva o somatório para o \textbf{melhor caso} (vetor já ordenado).
    \item Calcule as formas fechadas de ambos os casos e determine as complexidades.
\end{enumerate}

\begin{resposta}
\vspace{4cm}
\end{resposta}

\newpage
% ---------------------------------------------------------
\subsection*{Exercício 7 — Shell Sort em Java}

\begin{lstlisting}[style=javastyle]
void shellSort(int[] v) {
    int n = v.length;
    for (int gap = n / 2; gap > 0; gap /= 2) {       // gap: n/2, n/4, ...
        for (int i = gap; i < n; i++) {
            int tmp = v[i];
            int j = i;
            while (j >= gap && v[j - gap] > tmp) {   // comparacao
                v[j] = v[j - gap];
                j -= gap;
            }
            v[j] = tmp;
        }
    }
}
\end{lstlisting}

\begin{enumerate}[label=(\alph*)]
    \item Quantas vezes o laço externo (\texttt{gap}) executa em função de $n$? Expresse como somatório.
    \item O Shell Sort é mais difícil de analisar exatamente. Para o pior caso com a sequência de gaps $\lfloor n/2^k \rfloor$, a complexidade é $O(n^2)$. Tente montar um somatório que justifique isso.
    \item Por que o Shell Sort tende a ser mais rápido que o Insertion Sort na prática?
\end{enumerate}

\begin{resposta}
\vspace{4cm}
\end{resposta}

% =========================================================
\section*{Parte III — Desafio}
% =========================================================

% ---------------------------------------------------------
\subsection*{Exercício 8 — Comparação entre algoritmos}

Preencha a tabela abaixo com o somatório exato e a forma fechada para cada algoritmo, considerando o \textbf{pior caso}:

\bigskip

\begin{center}
\renewcommand{\arraystretch}{2.5}
\begin{tabular}{|l|c|c|c|}
\hline
\textbf{Algoritmo} & \textbf{Somatório} & \textbf{Forma Fechada} & $\Theta(\;)$ \\
\hline
Selection Sort & & & \\
\hline
Bubble Sort    & & & \\
\hline
Insertion Sort & & & \\
\hline
Shell Sort     & & & \\
\hline
\end{tabular}
\end{center}

\bigskip

% ---------------------------------------------------------
\subsection*{Exercício 9 — Código misterioso em C}

Dado o código abaixo, \textbf{sem olhar o nome da variável}, determine o somatório e identifique qual algoritmo de ordenação está sendo representado:

\begin{lstlisting}[style=cstyle]
void misterio(int v[], int n) {
    int i, j, x;
    for (i = 0; i < n - 1; i++) {
        x = i;
        for (j = i + 1; j < n; j++)
            if (v[j] < v[x]) x = j;
        int tmp = v[x]; v[x] = v[i]; v[i] = tmp;
    }
}
\end{lstlisting}

\begin{enumerate}[label=(\alph*)]
    \item Monte o somatório do número de comparações.
    \item Calcule a forma fechada.
    \item Qual algoritmo é este?
\end{enumerate}

\begin{resposta}
\vspace{3.5cm}
\end{resposta}

% ---------------------------------------------------------
\subsection*{Exercício 10 — Crie seu próprio código}

Escreva um trecho de código (em C ou Java) cujo número de operações seja descrito pelo somatório abaixo, e então verifique calculando a forma fechada:

\[
\sum_{i=1}^{n} \sum_{j=i}^{n} 1
\]

\begin{resposta}
\vspace{4cm}
\end{resposta}

\end{document}
